\documentclass[10pt,a4paper]{amsart}
\usepackage[margin=19mm]{geometry}
\usepackage{amsmath,amssymb,mathtools}
\usepackage[scr=rsfs,cal=euler]{mathalfa}
\usepackage{xfrac}
\usepackage[unicode,hidelinks,bookmarks=false]{hyperref}
\newcommand{\ie}{i.e.\ }
\newcommand{\cf}{cf.\ }
\newcommand{\resp}{respectively\ }
\newcommand{\Subsection}{\subsection}
\providecommand{\nobreakdash}{\nobreak\hbox{-}}
\setlength{\emergencystretch}{2em}
\multlinegap0pt
\usepackage{amssymb}
\usepackage{xcolor}
\newtheorem{thm}{Theorem}
\newtheorem{prop}{Proposition}
\newcommand{\CC}{\ensuremath{\mathbb{C}}}
\newcommand{\Cg}{\ensuremath{C_{\ol g}}}
\newcommand{\Hinf}{\ensuremath{H^\infty}}
\newcommand{\Hp}{\ensuremath{H^2_+}}
\newcommand{\OO}{\ensuremath{\mathcal{O}}}
\newcommand{\ol}[1]{\ensuremath{\overline{#1}}}
\title{Maximal functions, conjugations and multipliers between Toeplitz kernels}
\author{M. Cristina C\^amara, C. Carteiro, C. Diogo}
\date{Source: arXiv:2512.20406v1 (CC BY 4.0)}
\begin{document}
\maketitle

\noindent\emph{Source and licence.} Maximal functions, conjugations and multipliers between Toeplitz kernels. Version arXiv:2512.20406v1 (CC BY 4.0), distributed by arXiv under the Creative Commons Attribution 4.0 International licence, http://creativecommons.org/licenses/by/4.0/, as recorded in the arXiv licence metadata for this exact version and shown on the abstract page https://arxiv.org/abs/2512.20406v1. Full licence notice: \texttt{licenses/arxiv-2512.20406-CC-BY-4.0.txt}.\par\smallskip

\noindent\textit{Editorial scope.} Counted unit: the article's prop prop:maximal\_function\_prescribed\_inner\_factor (source0.tex lines 1204--1213). The article's complete original proof of it is delivered with it (source0.tex lines 1215--1226, one \texttt{proof} environment of 12 lines). No mathematics is written, repaired or re-derived: the statement and the proof are the article's own lines, in the article's own notation, and no retained line is substituted or re-flowed.  Retained uncounted as the source-local context the proof needs: lines 1055--1058.  Nothing was omitted from any retained span, so the package carries no omission record.  No retained line contains a citation, so the wrapper carries no bibliography and no bibliography entry.  No retained line contains a figure environment, an image-inclusion command or drawing code, so no image is delivered and the package declares no asset dependency.  Editorial formatting only, in addition: an AMS \texttt{amsart} preamble that restates the article's own declarations and macros used by the retained lines, namely the environments \texttt{thm}, \texttt{prop} on separate counters, together with the standard packages those lines need.  The retained environments print in this document as Theorem 1, Proposition 1 in order of appearance, whereas the article prints the same items with the numbering of its own full document.  The complete native source of the article as distributed in the arXiv e-print for this exact version is bundled unchanged as \texttt{sources/191565/source0.tex}.
\begin{thm}
\label{thm:max_func_charact_with_conj}
Let $f\in\Hp$; $f$ is a maximal function in $\ker T_g$ if and only if $\Cg f$ is outer in \Hp.
\end{thm}

\begin{prop}
\label{prop:maximal_function_prescribed_inner_factor}
For every inner function $\alpha$ such that $\alpha \prec I$, there exists a maximal function $f_\alpha^M$ in $\ker T_g$, with inner-outer factorisation $f_\alpha^M = \alpha\OO$, where $\OO\in\Hp$ is outer, given by
\begin{equation}
\label{eq:outer_factor_max_func_prescribed_inner_factor}
\OO = O (\lambda + I \ol \alpha ), \text{ with } \lambda \in\CC, |\lambda| = 1.
\end{equation}

We have then $f_\alpha^M \prec f^M$ and $\ker T_{\ol z \frac{\ol\OO}{\OO}} = \ker T_{\ol z \alpha \ol I \frac{\ol O}{O}} \supsetneq \ker T_{\ol z \frac{\ol O}{O}}$.
\end{prop}

\begin{proof}
For any $\lambda \in \CC$, $|\lambda|=1$, $\lambda + I \ol\alpha$ is outer in \Hinf, so $\OO$ is outer and
\begin{align*}
\Cg ( \alpha\OO) &= \Cg \left(\alpha O \left(\lambda + I \ol\alpha\right) \right) = \ol {gz} \ol O \left(\ol{ I + \lambda \alpha}\right) \\
&= \left(z I \frac{O}{\ol O}\right) \ol z \ol O \left(\ol{ I + \lambda \alpha}\right) = \ol\lambda O \left(\lambda + I \ol\alpha\right)
\end{align*}
where the right-hand side represents an outer function in \Hp.
Thus, by Theorem \ref{thm:max_func_charact_with_conj}, $f_\alpha^M = \alpha O$ is a maximal function in $\ker T_g$ such that $f_\alpha^M \prec f^M$ and the last inclusion follows from the equality
\[
\frac{\ol\OO}{\OO} = \frac{\ol O (\ol\lambda - \ol I \alpha)}{O (\lambda - I \ol\alpha)} = \frac{\ol O}{O} \ol\lambda \ol I \alpha.
\]
\end{proof}

\end{document}
